\section{What the correction is worth}
\label{sec:worth}

Sections \ref{sec:replication} and \ref{sec:mechanism} leave an obvious
question unanswered: how much \emph{should} FID have moved? Without that
number, ``the metrics did not respond'' is an observation rather than an
argument.

We answer it by applying the correction directly. Taking each model's own
samples and boosting the high band by \SI{0.41}{\decibel} --- the amount the
trained objective achieves --- is a pure post-process, so the resulting change
in FID isolates the value of the spectral correction from everything else
training might have altered.

\begin{table}[t]
\centering
\caption{The correction applied as a Fourier post-process, five seeds per arm.
The \SI{0}{\decibel} column reproduces the sweep's FID exactly, validating the
apparatus.}
\label{tab:worth}
\begin{tabular}{lrrr}
\toprule
& Spectral deficit & FID & FID after $+\SI{0.41}{\decibel}$ \\
\midrule
MSE & \SI{2.22}{\decibel} & $\fidMSE \pm 0.210$ & $\mathbf{\fidMSEboost \pm 0.202}$ \\
Frequency-aware & \SI{1.80}{\decibel} & $\fidAWWL \pm 0.344$ & $\fidAWWLboost \pm 0.335$ \\
\bottomrule
\end{tabular}
\end{table}

Two independent measurements agree that the objective corrects
$\correction \pm 0.18$~\si{\decibel}: the difference of deficits in
Table~\ref{tab:worth}, and the band analysis of \S\ref{sec:mechanism} obtained
by a different route.

Applying that correction post hoc is worth $\boostValue \pm \boostCI$~FID. The
accounting follows immediately.

\begin{table}[t]
\centering
\caption{The objective falls \shortfall~FID short of what its own spectral
correction predicts.}
\label{tab:accounting}
\begin{tabular}{lr}
\toprule
& FID \\
\midrule
MSE baseline & \fidMSE \\
Predicted for the objective from its own spectral correction & \fidMSEboost \\
Objective, measured & \fidAWWL \\
\midrule
\textbf{Unexplained penalty} & $\mathbf{+\shortfall}$ \\
\bottomrule
\end{tabular}
\end{table}

The objective buys $0.35$~FID of spectral improvement and pays roughly
$1.1$~FID for it elsewhere. Two framings of the same fact are worth stating
separately:

\begin{itemize}
  \item MSE plus a twenty-line Fourier post-process (\fidMSEboost) outperforms
    the trained objective (\fidAWWL) by \shortfall~FID, at no training cost.
  \item The objective plus the same post-process (\fidAWWLboost) is
    \emph{still} worse than plain MSE (\fidMSE). The correction cannot recover
    what obtaining it costs.
\end{itemize}

\paragraph{Measurement hazard: calibrate away from the minimum.}
We first calibrated FID's sensitivity on real images --- attenuating one half
of a real set and scoring it against the other. That measures sensitivity at
the \emph{minimum} of the distance, where the first derivative vanishes and
perturbations are penalised quadratically. It implied a \SI{0.41}{\decibel}
correction was worth $0.010$~FID and would require some 17\,500 seeds to
detect. Measured at the operating point, where the derivative is non-zero, the
same correction is worth $\boostValue$ --- thirtyfold larger. Sensitivity
calibrations for a generative metric must be performed away from the minimum;
we drew, and retract, the opposite conclusion from the wrong calibration.
